\section{Introduction}
\label{sec:intro}

% ── Mark scheme: 10 marks, ~half a page ──────────────────────────────────────

Game software normally runs entirely on a PC.
The PC decides whether a move is legal, tracks the board, and detects the winner.
This design has a serious weakness: a client that sends moves in the wrong order, or
claims a square is free when it is not, could corrupt the game.
The project brief~\cite{brief} addresses this by reversing the roles.
The STM32 microcontroller is the sole referee.
It stores the board, validates every action and returns the authoritative result.
The PC client only draws the board and forwards what the player clicks.
The central design invariant enforced by the server is that a rejected request
leaves the game state strictly unchanged---the board after an invalid move is
identical to the board before it.

We implemented \textbf{Tic-Tac-Toe} (also called noughts and crosses), the baseline
game described in the brief.
Tic-Tac-Toe has a small, well-defined state (nine cells, two players, eight winning
lines) that fits comfortably in the STM32's memory and can be exhaustively tested
in under a minute.
For the graphical user interface (GUI) we chose Python's built-in \textbf{Tkinter}
library~\cite{guilib}.
Tkinter requires no additional installation, runs on Windows without extra drivers,
and handles user events and screen updates in a single thread using the Tk event loop,
which removes the need for explicit threads when waiting for the serial reply.

The system is divided into three layers.
The embedded firmware on the STM32 handles all rule enforcement and game logic.
A thin serial driver on the PC opens the COM port and exchanges newline-terminated
ASCII frames at 115\,200 baud.
The Tkinter GUI sits on top of the driver and deals only with drawing and user input.
This layering means the game rules can be tested independently of the screen by
sending raw commands over the serial port, which is exactly how our automated test
harness works.

\textbf{Task allocation.}
Nemukula Vhugala (Partner~A) was responsible for the embedded firmware: the
interrupt-driven UART receiver and circular ring buffer (\texttt{protocol.c}),
the game finite-state machine (FSM) and rule engine (\texttt{game.c}), and the
hardware integration---button debouncing and LED status indicators
(\texttt{hardware.c}).
Munyai Umusindisi Cindy (Partner~B) was responsible for the PC-side software:
the serial communication driver (\texttt{serial\_client.py}), the Tkinter GUI
(\texttt{gui.py}), and the automated hardware test harness
(\texttt{test\_day3\_stm32.py}).
Both partners contributed to the protocol specification, system integration,
and this report.

\textbf{Structure.}
Section~\ref{sec:protocol} describes the serial link, frame format, and full
command set.
Section~\ref{sec:design} presents the server FSM, GUI flowchart, and circuit diagram.
Section~\ref{sec:impl} highlights the most challenging pieces of implementation.
Section~\ref{sec:authority} argues that the server enforces all game rules
and cannot be bypassed by a misbehaving client.
Section~\ref{sec:robust} explains how the server handles malformed or unexpected input.
Section~\ref{sec:valid} presents the automated test results on the real STM32 board.
Section~\ref{sec:conclusion} summarises the outcome and suggests improvements.
